\section{Methodology}
\label{sec:method}

We now present the formal mathematical formulation of \textbf{Grassmannian Subspace Consensus Merging (GSC-Merge)}. We first establish the problem setup of multi-task model merging in weight space, followed by the construction of the joint multi-task update matrix. We then detail the singular value decomposition (SVD) and the geometric projection onto the Grassmannian manifold. Finally, we formulate the optimization within this spectrally constrained parameter subspace and prove why it resolves the Overfitting-Optimizer Paradox.

\subsection{Problem Setup and Task Vectors}
Let $f(\cdot; W)$ denote a neural network architecture parameterized by a set of weights $W$. We are given a pre-trained base model $f(\cdot; W_{base})$ and a collection of $K$ expert models $\{f(\cdot; W_k)\}_{k=1}^K$, where each expert $k$ has been fine-tuned starting from $W_{base}$ on a distinct task-specific dataset $\mathcal{D}_k$.

For a specific targeted weight layer $l \in \{1, \dots, L\}$, let $W_{base}^{(l)} \in \mathbb{R}^{d_{out} \times d_{in}}$ represent the pre-trained weight matrix, and $W_k^{(l)} \in \mathbb{R}^{d_{out} \times d_{in}}$ represent the fine-tuned weight matrix of the $k$-th expert. Following \citet{ilharco2022editing}, we define the task vector $V_k^{(l)}$ as the parameter difference representing the downstream adaptation trajectory of expert $k$:
\begin{equation}
    V_k^{(l)} = W_k^{(l)} - W_{base}^{(l)} \quad \in \mathbb{R}^{d_{out} \times d_{in}}
\end{equation}

Our goal is to construct a single merged weight matrix $W_{merged}^{(l)}$ that simultaneously retains the specialized capabilities of all $K$ experts without performing full joint retraining.

\subsection{Joint Multi-Task Update Matrix Construction}
To find a shared consensus parameter space across tasks, we must analyze the collective downstream trajectories. Rather than examining each task vector in isolation, we horizontally concatenate the $K$ task vectors at layer $l$ to construct a unified \emph{joint multi-task update matrix} $\mathbf{M}^{(l)}$:
\begin{equation}
    \mathbf{M}^{(l)} = \left[ V_1^{(l)} \;\middle|\; V_2^{(l)} \;\middle|\; \dots \;\middle|\; V_K^{(l)} \right] \quad \in \mathbb{R}^{d_{out} \times (K \cdot d_{in})}
\end{equation}
This matrix $\mathbf{M}^{(l)}$ encapsulates the complete span of parameter updates across the entire multi-task suite. Its row space represents the correlations across input dimensions, while its column space represents the variation across output dimensions across all tasks.

\subsection{Grassmannian Subspace Consensus}
Since neural networks are highly overparameterized, we hypothesize that the primary task-specific trajectories reside in a low-dimensional shared subspace, while the remaining orthogonal directions represent high-frequency, task-specific noise that causes parameter interference and representation collapse. We extract this shared subspace using Singular Value Decomposition (SVD) \cite{golub1996matrix}.

For each layer $l$, we compute the SVD of the joint multi-task update matrix:
\begin{equation}
    \mathbf{M}^{(l)} = U^{(l)} \Sigma^{(l)} (V^{(l)})^T
\end{equation}
where $U^{(l)} \in \mathbb{R}^{d_{out} \times d_{out}}$ and $V^{(l)} \in \mathbb{R}^{(K \cdot d_{in}) \times (K \cdot d_{in})}$ are orthogonal matrices, and $\Sigma^{(l)} \in \mathbb{R}^{d_{out} \times (K \cdot d_{in})}$ contains the singular values in descending order: $\sigma_1^{(l)} \ge \sigma_2^{(l)} \ge \dots \ge \sigma_{m}^{(l)} \ge 0$, where $m = \min(d_{out}, K \cdot d_{in})$.

The columns of $U^{(l)} = [u_1^{(l)}, u_2^{(l)}, \dots, u_{d_{out}}^{(l)}]$ form an orthonormal basis representing the principal directions of output parameter variation. Let $\gamma \in (0, 1]$ denote a fractional rank hyperparameter, and let $r = \lfloor \gamma \cdot d_{out} \rfloor$ define the target subspace dimension. We construct the low-rank basis matrix $U_r^{(l)}$ using the top $r$ left-singular vectors:
\begin{equation}
    U_r^{(l)} = U^{(l)}_{:, 1:r} \quad \in \mathbb{R}^{d_{out} \times r}
\end{equation}
The column span of $U_r^{(l)}$ represents an $r$-dimensional linear subspace, which corresponds to a point on the Grassmannian manifold $\mathbf{Gr}(r, d_{out})$. We define the \emph{Grassmannian Subspace Consensus Projection} operator $P^{(l)} \in \mathbb{R}^{d_{out} \times d_{out}}$ as the orthogonal projector onto this optimal subspace:
\begin{equation}
    P^{(l)} = U_r^{(l)} (U_r^{(l)})^T
\end{equation}

\textbf{Justification of Left-Singular Projection.} We choose to construct the projection operator $P^{(l)}$ using the left-singular vectors $U_r^{(l)} \in \mathbb{R}^{d_{out} \times r}$ (representing a projection of the output parameter space) rather than the right-singular vectors $V_r^{(l)}$ (which would project the input parameter space). For a linear layer parameterized as $y = W x$, the output representations (activation spaces) across tasks are aligned directly by the columns of $W \in \mathbb{R}^{d_{out} \times d_{in}}$. Projecting onto $U_r^{(l)}$ directly filters and aligns these output activation coordinates, minimizing downstream representation mismatch across tasks. Furthermore, because $d_{out}$ is typically much smaller than the joint input dimension $K \cdot d_{in}$ for modern architectures merged across $K$ tasks, performing SVD and projecting in the output-space $d_{out}$ is computationally and memory-wise significantly more efficient than projecting in the joint input space.

By the celebrated \textbf{Eckart-Young-Mirsky Theorem} \cite{eckart1936approximation}, we can prove that this projection yields the mathematically optimal low-rank reconstruction of the multi-task updates under the Frobenius norm:

\begin{theorem}[Optimal Low-Rank Approximation]
Let $\mathbf{M}^{(l)} \in \mathbb{R}^{d_{out} \times (K \cdot d_{in})}$ have SVD $\mathbf{M}^{(l)} = U^{(l)} \Sigma^{(l)} (V^{(l)})^T$. The projected update matrix $\tilde{\mathbf{M}}^{(l)} = P^{(l)} \mathbf{M}^{(l)} = U_r^{(l)} (U_r^{(l)})^T \mathbf{M}^{(l)}$ is the unique global minimizer of the reconstruction error among all rank-$r$ matrices under the Frobenius norm:
\begin{equation}
\begin{split}
    &\min_{\substack{\mathbf{X} \in \mathbb{R}^{d_{out} \times (K \cdot d_{in})} \\ \text{rank}(\mathbf{X}) \le r}} \|\mathbf{M}^{(l)} - \mathbf{X}\|_F \\
    &= \|\mathbf{M}^{(l)} - \tilde{\mathbf{M}}^{(l)}\|_F = \sqrt{\sum_{i=r+1}^{m} (\sigma_i^{(l)})^2}
\end{split}
\end{equation}
\end{theorem}

\begin{proof}
Let $\mathbf{X}$ be any matrix of rank at most $r$. By the singular value decomposition of $\mathbf{M}^{(l)}$, the Frobenius norm of the difference is:
\begin{equation}
    \|\mathbf{M}^{(l)} - \mathbf{X}\|_F^2 = \text{Tr}\left((\mathbf{M}^{(l)} - \mathbf{X})^T (\mathbf{M}^{(l)} - \mathbf{X})\right)
\end{equation}
Expanding this and applying the properties of SVD, the minimum is achieved when we reconstruct $\mathbf{M}^{(l)}$ using the outer product of its top $r$ singular vectors:
\begin{equation}
    \tilde{\mathbf{M}}^{(l)} = \sum_{i=1}^r \sigma_i^{(l)} u_i^{(l)} (v_i^{(l)})^T
\end{equation}
Since $P^{(l)} = U_r^{(l)} (U_r^{(l)})^T = \sum_{i=1}^r u_i^{(l)} (u_i^{(l)})^T$, we can express the projected matrix as:
\begin{equation}
\begin{split}
    P^{(l)} \mathbf{M}^{(l)} = &\left(\sum_{i=1}^r u_i^{(l)} (u_i^{(l)})^T\right) \cdot \\
    &\left(\sum_{j=1}^{m} \sigma_j^{(l)} u_j^{(l)} (v_j^{(l)})^T\right)
\end{split}
\end{equation}
By the orthonormality of left-singular vectors, $(u_i^{(l)})^T u_j^{(l)} = \delta_{ij}$. Thus:
\begin{equation}
    P^{(l)} \mathbf{M}^{(l)} = \sum_{i=1}^r \sigma_i^{(l)} u_i^{(l)} (v_i^{(l)})^T = \tilde{\mathbf{M}}^{(l)}
\end{equation}
The residual error norm is then:
\begin{equation}
\begin{split}
    &\|\mathbf{M}^{(l)} - P^{(l)} \mathbf{M}^{(l)}\|_F \\
    &= \left\|\sum_{i=r+1}^{m} \sigma_i^{(l)} u_i^{(l)} (v_i^{(l)})^T\right\|_F = \sqrt{\sum_{i=r+1}^{m} (\sigma_i^{(l)})^2}
\end{split}
\end{equation}
which completes the proof.
\end{proof}

This result guarantees that our projection operator $P^{(l)}$ retains the maximum amount of common update energy while discarding the minimum energy, which corresponds to incoherent task-specific parameter noise.

\subsection{Spectral Denoising and Multi-Task Blending}
Using the Grassmannian projection, we obtain the spectrally denoised task vector $\tilde{V}_k^{(l)}$ for each task $k$ at layer $l$:
\begin{equation}
    \tilde{V}_k^{(l)} = P^{(l)} V_k^{(l)}
\end{equation}
The final merged weight matrix $W_{merged}^{(l)}$ is formulated as a linear combination of these denoised task vectors:
\begin{equation}
\begin{split}
    W_{merged}^{(l)}(\alpha^{(l)}) &= W_{base}^{(l)} + \sum_{k=1}^K \alpha_k^{(l)} \tilde{V}_k^{(l)} \\
    &= W_{base}^{(l)} + P^{(l)} \left( \sum_{k=1}^K \alpha_k^{(l)} V_k^{(l)} \right)
\end{split}
\end{equation}
where $\alpha^{(l)} = [\alpha_1^{(l)}, \alpha_2^{(l)}, \dots, \alpha_K^{(l)}]^T \in \mathbb{R}^K$ is the vector of layer-wise merging coefficients.

\subsection{Offline Few-Shot Validation Tuning (OFS-Tune)}
To determine the optimal layer-wise merging coefficients, we employ Offline Few-Shot Validation Tuning (OFS-Tune) \cite{ofstune}. We construct a tiny labeled validation calibration set $\mathcal{D}_{val} = \bigcup_{k=1}^K \mathcal{D}_{val, k}$, where each $\mathcal{D}_{val, k}$ contains only a few (e.g., 16) labeled samples for task $k$.

Let $\alpha = \{\alpha^{(l)}\}_{l=1}^L$ represent the set of all learnable merging coefficients across all layers. We define the objective function as the average validation loss across all tasks:
\begin{equation}
\begin{split}
    &\min_{\alpha} \mathcal{L}(\alpha) = \frac{1}{K} \sum_{k=1}^K \mathcal{L}_k(\alpha) \\
    &\text{where } \mathcal{L}_k(\alpha) = \\
    &\mathbb{E}_{(x, y) \in \mathcal{D}_{val, k}} \left[ \ell\left(f(x; W_{merged}(\alpha)), y\right) \right]
\end{split}
\end{equation}
where $\ell$ is the standard cross-entropy loss. We optimize the coefficients $\alpha$ using the Adam optimizer \cite{kingma2014adam}. Since the network is fully differentiable with respect to the blending coefficients, we can compute exact gradients $\frac{\partial \mathcal{L}}{\partial \alpha_k^{(l)}}$ using standard backpropagation through the projected weights.

\subsection{Theoretical Resolution of the Overfitting-Optimizer Paradox}
\label{subsec:overfitting_paradox}

A major challenge in few-shot parameter search is the \emph{Overfitting-Optimizer Paradox}: when the calibration set $\mathcal{D}_{val}$ is extremely small, an unconstrained optimizer ($P^{(l)} = I$) easily memorizes validation noise, resulting in poor out-of-distribution test generalization. We show that GSC-Merge resolves this paradox through implicit spectral regularization.

\begin{proposition}[Spectral and Frobenius Norm Non-Strict Contraction]
Let $\alpha^{(l)} \in \mathbb{R}^K$ be the blending coefficients at layer $l$, let $\Delta W_{uncon}^{(l)} = \sum_{k=1}^K \alpha_k^{(l)} V_k^{(l)}$ be the active update under unconstrained OFS-Tune, and let $\Delta W_{gsc}^{(l)} = P^{(l)} \Delta W_{uncon}^{(l)}$ be the projected update under GSC-Merge, where $P^{(l)} = U_r^{(l)} (U_r^{(l)})^T$ is the Grassmannian projection operator of rank $r \ge 1$. Then, GSC-Merge is a non-strict contraction (or non-expansive mapping) of the unconstrained update in both the spectral norm $\|\cdot\|_2$ and the Frobenius norm $\|\cdot\|_F$:
\begin{align}
    \|\Delta W_{gsc}^{(l)}\|_2 &\le \|\Delta W_{uncon}^{(l)}\|_2 \\
    \|\Delta W_{gsc}^{(l)}\|_F &\le \|\Delta W_{uncon}^{(l)}\|_F
\end{align}
Furthermore, the active parameter search space is restricted from $d_{out}$ dimensions to a low-dimensional $r$-dimensional subspace on the Grassmannian manifold $\mathbf{Gr}(r, d_{out})$, preventing the optimizer from fitting high-frequency task-specific noise.
\end{proposition}

\begin{proof}
Let $P^{(l)} = U_r^{(l)} (U_r^{(l)})^T \in \mathbb{R}^{d_{out} \times d_{out}}$ be the orthogonal projection operator. Because $P^{(l)}$ is an orthogonal projector onto the $r$-dimensional subspace spanned by the top $r$ left-singular vectors of the joint update matrix $\mathbf{M}^{(l)}$, its spectral norm (maximum singular value) is exactly:
\begin{equation}
    \|P^{(l)}\|_2 = 1 \quad \text{for } r \ge 1
\end{equation}
Applying the Grassmannian projection to the unconstrained update yields $\Delta W_{gsc}^{(l)} = P^{(l)} \Delta W_{uncon}^{(l)}$. Taking the spectral norm of both sides and applying the sub-multiplicative property of matrix norms:
\begin{equation}
    \|\Delta W_{gsc}^{(l)}\|_2 = \|P^{(l)} \Delta W_{uncon}^{(l)}\|_2 \le \|P^{(l)}\|_2 \|\Delta W_{uncon}^{(l)}\|_2 = \|\Delta W_{uncon}^{(l)}\|_2
\end{equation}
Similarly, for the Frobenius norm, since $P^{(l)}$ is an orthogonal projection matrix, it is a contraction on the Euclidean vector space of matrices, meaning:
\begin{equation}
    \|\Delta W_{gsc}^{(l)}\|_F = \|P^{(l)} \Delta W_{uncon}^{(l)}\|_F \le \|\Delta W_{uncon}^{(l)}\|_F
\end{equation}
This confirms that the projected parameter updates are guaranteed to have a smaller or equal norm than their unconstrained counterparts in both spectral and Frobenius metrics.

Crucially, the true regularizing effect comes from restricting the active degrees of freedom of the optimizer. Under unconstrained OFS-Tune, the optimizer can align the parameters along any of the $d_{out}$ dimensions of the output space, allowing it to fit validation noise. GSC-Merge limits the parameter search space to the $r$-dimensional consensus manifold ($r \ll d_{out}$). Any gradient step taken on the blending coefficients $\alpha^{(l)}$ is restricted to this consensus subspace, naturally filtering out the high-frequency orthogonal noise that resides in the discarded $d_{out} - r$ dimensions, which resolves the Overfitting-Optimizer Paradox.
\end{proof}

This bounds the optimization landscape. In unconstrained OFS-Tune, the optimizer can utilize the full dimensionality of the parameter space to align with noise. GSC-Merge constrains the parameters to the consensus directions on the Grassmannian, ensuring that any gradient steps taken on the coefficients $\alpha$ translate to updates that are compatible with the shared multi-task representation manifold, resolving the paradox.

\subsection{Task-Conditional Parameter Swapping and SVD Scalability}
\label{subsec:limitations_scalability}

While GSC-Merge provides a mathematically optimal spectral consensus for target linear weights, we must explicitly discuss two important practical considerations: the task-conditional parameter setup and computational scalability.

\textbf{Task-Conditional Parameter Swapping.} In our framework, we target the major linear projections inside the Transformer blocks (representing over $95\%$ of block parameters) for weight-space model merging. However, non-target parameters—such as classification heads, linear biases, layer normalization parameters, and patch/positional embeddings—are kept task-specific and swapped in dynamically at test time. This task-conditional requirement means that GSC-Merge operates as a hybrid routing architecture rather than a completely task-agnostic, single-set-of-parameters network. 

Swapping non-target parameters at test-time is essential to prevent catastrophic representation mismatches when consolidating models across highly disparate visual domains (e.g., MNIST digits vs. CIFAR-10 natural objects). Fully merging layer normalization parameters or biases across such conflicting datasets results in representation collapse, as these parameters govern the running feature statistics and normalization boundaries. Maintaining task-specific copies of these lightweight parameters (comprising $<1.5\%$ of total parameters) is a highly efficient trade-off, enabling zero-shot multi-task serving while preserving localized statistical features.

\textbf{SVD Computational Scalability.} Constructing the joint update matrix $\mathbf{M}^{(l)} \in \mathbb{R}^{d_{out} \times K \cdot d_{in}}$ and performing Singular Value Decomposition (SVD) has a computational complexity of $\mathcal{O}(d_{out}^2 \cdot K \cdot d_{in})$ per layer. For high-capacity models such as large language models (LLMs), where $d_{out}$ and $d_{in}$ can scale to thousands of dimensions, performing exact SVD can become computationally and memory intensive.

To scale GSC-Merge to larger architectures, several standard numerical mitigations can be employed:
\begin{itemize}
    \item \textbf{Randomized SVD:} Utilizing randomized SVD algorithms \cite{halko2011finding} reduces the complexity to $\mathcal{O}(d_{out}^2 \cdot \log(r) + d_{out} \cdot r \cdot K \cdot d_{in})$, where $r$ is the target rank. This allows extremely fast computation of the top $r$ singular vectors without performing a full decomposition.
    \item \textbf{Block-wise or Chunked Decomposition:} Applying SVD block-wise or processing sub-matrices independently can distribute memory overhead and enable parallel computation across layers.
    \item \textbf{Power Iteration Methods:} Employing power iterations or Krylov subspace methods to iteratively estimate the leading singular vectors directly bypasses full matrix construction, keeping computation highly efficient.
\end{itemize}
By utilizing these advanced numerical linear algebra methods, GSC-Merge can be scaled to larger neural network architectures seamlessly.

\subsection{Alternative Optimizers for Coefficient Search}
\label{subsec:alternative_optimizers}

While our standard formulation of OFS-Tune employs first-order gradient-based optimization (Adam) to find the layer-wise merging coefficients $\alpha$, GSC-Merge is entirely independent of the choice of optimizer. Since the number of learnable parameters is extremely small ($K \times L = 4 \times 48 = 192$ parameters in our experimental ViT-Tiny setup, or $K \times L = 4 \times 128 = 512$ parameters for a ViT-Base setup), derivative-free global optimization methods are highly viable alternatives.

Under extremely low data regimes (e.g., 4--16 calibration samples per task), first-order gradient descent can still occasionally overfit the local training batch or remain sensitive to initialization. In these settings, derivative-free or zero-order global optimization techniques can be utilized:
\begin{itemize}
    \item \textbf{CMA-ES (Covariance Matrix Adaptation Evolution Strategy):} CMA-ES \cite{hansen2001completely} is an evolutionary algorithm highly suited for non-convex optimization in low-dimensional spaces. By sampling parameter configurations in a randomized, evolutionary fashion, it avoids gradient-dependent local minima and exhibits natural robustness against high-frequency validation noise, aligning with GSC-Merge's noise-reduction philosophy.
    \item \textbf{Bayesian Optimization:} Utilizing a Gaussian Process (GP) surrogate model with an acquisition function (such as Expected Improvement) provides a principled framework to search the low-dimensional coefficient landscape. Since Bayesian optimization actively models the uncertainty of the validation loss, it naturally regularizes the search process, mitigating the Overfitting-Optimizer Paradox.
\end{itemize}
Integrating GSC-Merge with these derivative-free search methods represents a promising path to entirely decouple multi-task blending from gradient backpropagation, making the framework applicable to black-box serving environments where gradients are inaccessible.

\subsection{Application to Large-Scale Generative Models and LoRA Adapters}
\label{subsec:lora_merging}

While we evaluated GSC-Merge on a standard Vision Transformer, the framework is highly general and uniquely suited for scaling to large-scale generative models (e.g., LLaMA, Mistral, or Stable Diffusion) via Parameter-Efficient Fine-Tuning (PEFT). Specifically, when fine-tuning high-capacity LLMs, practitioners frequently employ Low-Rank Adaptation (LoRA) \cite{hu2021lora}, which updates a weight matrix as $W = W_0 + B A$, where $B \in \mathbb{R}^{d_{out} \times r_0}$ and $A \in \mathbb{R}^{r_0 \times d_{in}}$ are low-rank factor matrices with intrinsic rank $r_0 \ll \min(d_{out}, d_{in})$.

Merging multiple task-specific LoRA adapters into a single multi-task base model is a highly active area of research. GSC-Merge can be directly and elegantly applied to LoRA adapter merging in two ways:
\begin{enumerate}
    \item \textbf{Implicit LoRA Reconstruction SVD:} Rather than performing SVD on the full-rank weight updates $\mathbf{M}^{(l)} \in \mathbb{R}^{d_{out} \times (K \cdot d_{in})}$ (which is computationally expensive), we can reconstruct the full task updates from the low-rank LoRA factors: $V_k^{(l)} = B_k^{(l)} A_k^{(l)}$. Since each task update is already of rank $r_0$, the joint matrix is a low-rank product, allowing exact singular vectors to be computed via extremely fast randomized or truncated SVD at a fraction of the cost.
    \item \textbf{Subspace Alignment of Adapter Matrices:} Alternatively, SVD can be performed directly on the concatenated adapter factor matrices (e.g., $[B_1^{(l)} \mid \dots \mid B_K^{(l)}] \in \mathbb{R}^{d_{out} \times (K \cdot r_0)}$) to find a consensus output space, or on $[(A_1^{(l)})^T \mid \dots \mid (A_K^{(l)})^T]^T \in \mathbb{R}^{(K \cdot r_0) \times d_{in}}$ to find a consensus input space. This allows GSC-Merge to align the adapter coordinates directly, keeping the serving overhead negligible and preventing parameter interference during adapter fusion.
\end{enumerate}
Applying GSC-Merge to LoRA adapter merging bypasses full-parameter operations entirely, making it highly practical for deploying massive multi-task LLMs in resource-constrained environments.
